\documentclass[conference]{IEEEtran}
\IEEEoverridecommandlockouts
% The preceding line is only needed to identify funding in the first footnote. If that is unneeded, please comment it out.
\usepackage{cite}
\usepackage{amsmath,amssymb,amsfonts}
\usepackage{algorithmic}
\usepackage{graphicx}
\usepackage{textcomp}
\usepackage{xcolor}
\def\BibTeX{{\rm B\kern-.05em{\sc i\kern-.025em b}\kern-.08em
    T\kern-.1667em\lower.7ex\hbox{E}\kern-.125emX}}
\begin{document}

% The following line adds an experimental way to comment code out
% Thanks to https://tex.stackexchange.com/a/87305
\long\def\/*#1*/{}

\title{Cuckoo Hashing with Bounded Displacement for Exact-Match Flow Tables\\
\/*{\footnotesize \textsuperscript{*}Note: Sub-titles are not captured in Xplore and
should not be used}
\thanks{Identify applicable funding agency here. If none, delete this.}*/
}

\author{\IEEEauthorblockN{Soumyajit Kolay}
\IEEEauthorblockA{\textit{School of Computer Engineering} \\
\textit{KIIT Deemed University}\\
Bhubaneswar, India \\
2650026@kiit.ac.in}
}

\maketitle

\begin{abstract}
THE ABSTRACT OF THIS DOCUMENT IS YET TO BE UPDATED. I am learning \LaTeX as a part of this research.
\end{abstract}

\begin{IEEEkeywords}
SDN, hashtables, d-ary cuckoo, bloom filter, flow tables
\end{IEEEkeywords}

\section{Introduction}
The dictionary abstraction---supporting insertion, deletion, and membership queries over a dynamic key set---sits at the core of nearly every high-throughput networking data plane, and nowhere is its performance more scrutinized than in the exact-match flow table of a software-defined switch. Every packet traversing the data plane must be matched against this table at line rate, so a single slow lookup anywhere in the tail of the latency distribution translates directly into queuing delay or packet loss. Classical collision-resolution schemes such as chained hashing and open addressing with linear probing offer excellent average-case behavior but degrade to logarithmic or, in adversarial cases, linear worst-case lookup cost as the table fills---a property that is simply incompatible with a hardware or software pipeline built around a fixed per-packet time budget.

Cuckoo hashing was introduced by Pagh and Rodler as a direct answer to this problem~\cite{pagh2004cuckoo}. By allowing each key two candidate locations across two tables and permitting keys already resident in the table to be displaced (``kicked'') to their alternate location on collision, the scheme guarantees exactly two probes for every lookup---worst case, not merely expected. This property alone made cuckoo hashing an immediate candidate for line-rate dictionaries, and it has since been adopted widely in networking contexts precisely for this reason. Yet Pagh and Rodler's original construction carries a structural cost that later work has spent almost two decades trying to manage: when a chain of displacements returns to a previously visited cell, the insertion procedure has entered a cycle, and the only recourse in the original scheme is to abandon the current hash functions entirely and rehash the full table---an $O(n)$ operation that, however rare, is catastrophic if it occurs during a latency-sensitive insertion at line rate.

The subsequent literature can be read as a sustained effort to defer, shrink, or altogether avoid this full-table rehash. Kirsch, Mitzenmacher, and Wieder showed that reserving a small constant-sized side table---a stash---for the rare items that fail to settle after a bounded number of displacements reduces the probability of ever needing a full rehash from polynomial to a much smaller \$\textbackslash{}Theta(n\^{}\{-s-1\})\$ for a stash of size \$s\$, with a stash of only three or four entries already yielding a dramatic practical improvement\~{}\textbackslash{}cite\{kirsch2009stash\}. Kutzelnigg subsequently tightened this analysis, showing that the same failure-probability bounds hold even without the extra verification step Kirsch et al.'s original proof required, simplifying the practical insertion routine\~{}\textbackslash{}cite\{kutzelnigg2010stash\}. Aum\textbackslash{}"uller, Dietzfelbinger, and Woelfel then relaxed the hash-function assumptions underlying the stash guarantee, proving that simple, explicit 2-independent hash families suffice to preserve the \$O(n\^{}\{-s-1\})\$ failure bound and \$O(s)\$ query time, removing the reliance on fully random hash functions that had limited the scheme's practicality\~{}\textbackslash{}cite\{aumuller2014explicit\}. On the load-capacity side, Fountoulakis and Panagiotou resolved a long-open question by deriving the exact load thresholds for \$d\$-ary cuckoo hashing with \$d \textbackslash{}geq 3\$ choices per key---for instance, showing that three choices already permit roughly 92\textbackslash{}\% table utilization---giving practitioners, for the first time, precise rather than merely empirical guidance on how sparse a \$d\$-ary table must be kept\~{}\textbackslash{}cite\{fountoulakis2012sharp\}.

A parallel line of work adapted cuckoo hashing specifically to the networking substrate rather than treating it as a general-purpose dictionary. Pontarelli et al. proposed a parallel \$d\$-pipeline implementation targeting ASIC and FPGA memories accessed concurrently, reporting throughput gains of up to 60\textbackslash{}\% over prior parallel cuckoo designs for exact-match lookup\~{}\textbackslash{}cite\{pontarelli2016parallel\}, and later introduced Cuckoo Cache, which exploits the skewed popularity distribution of real network flows by biasing frequently accessed items toward the first-probed table, yielding substantial query-throughput improvements on real traffic traces without altering the underlying cuckoo guarantees\~{}\textbackslash{}cite\{pontarelli2016cuckoocache\}. Le Scouarnec pushed cache-efficiency further on commodity CPUs with Cuckoo++, an implementation tailored to Intel Xeon cache-line behavior that outperformed the widely used DPDK hash table by 45--70\textbackslash{}\% on lookups while leaving the classical theoretical analysis of cuckoo hashing untouched\~{}\textbackslash{}cite\{lescouarnec2018cuckoopp\}. Minaud and Papamanthou revisited and corrected a subtle error in the original multi-item-per-bucket stash analysis of Kirsch et al., establishing the true asymptotic failure probability as \$\textbackslash{}Theta(n\^{}\{-d-s\})\$ rather than the previously claimed bound\~{}\textbackslash{}cite\{minaud2023generalized\}---a correction of direct consequence to any system, like the one proposed here, that relies on stash-failure guarantees for its worst-case claims. Noble subsequently closed a longstanding gap in the theory by deriving the first explicit, closed-form failure-probability bounds for cuckoo hashing with a stash of arbitrary---not merely constant---size, removing the asymptotic hedging present in all prior stash analyses\~{}\textbackslash{}cite\{noble2021explicit\}.

Most recently, Kuszmaul and Mitzenmacher introduced bubble-up cuckoo hashing, a \$d\$-ary insertion policy that achieves an essentially optimal number of hash locations while guaranteeing \$O(\textbackslash{}delta\^{}\{-1\})\$ expected insertion time at any load factor \$1-\textbackslash{}delta\$ and \$O(1)\$ expected positive-query time independent of \$d\$---the strongest theoretical insertion-time guarantee for \$d\$-ary cuckoo hashing to date\~{}\textbackslash{}cite\{kuszmaul2025bubbling\}. On the hardware side, Meg\textbackslash{}'ias et al. designed an FPGA implementation that exploits parallel insertion iterations to cut the number of reallocation rounds by more than 60\textbackslash{}\% at loads beyond 95\textbackslash{}\%, targeting exactly the real-time communications-network setting this work is concerned with\~{}\textbackslash{}cite\{megias2025minimizing\}. Closest of all to the present application, Xiong et al. proposed TECache, a nearly conflict-free hashing scheme for TCAM-backed SDN flow tables that uses three candidate locations with co-directional kicking to sustain high cache-hit rates and substantial energy savings under real backbone traffic\~{}\textbackslash{}cite\{xiong2025tecache\}---a system that, notably, optimizes for hit rate and power rather than for a bound on worst-case insertion cost. Tripathi et al. most recently proposed a new set of performance indicators for sequential cuckoo hashing variants, aimed at giving practitioners richer evaluation vocabulary than raw throughput alone\~{}\textbackslash{}cite\{tripathi2026performance\}.

The stash-based line of theoretical work\~{}\textbackslash{}cite\{kirsch2009stash, kutzelnigg2010stash, aumuller2014explicit, minaud2023generalized, noble2021explicit\} establishes precisely \textbackslash{}emph\{how likely\} a full rehash is for a given stash size, but every one of these analyses assumes the classical insertion procedure: displacement continues, essentially without an explicit external bound, until either an empty slot is found or a cycle is detected, at which point the stash is invoked. None of them treat the length of the eviction chain itself as a tunable, deliberately constrained parameter---the stash is a safety net for whatever the unbounded process eventually fails to place, not a fallback triggered \textbackslash{}emph\{early\} by design. Conversely, the networking-systems line of work\~{}\textbackslash{}cite\{pontarelli2016parallel, pontarelli2016cuckoocache, lescouarnec2018cuckoopp, xiong2025tecache\} optimizes throughput, cache efficiency, or energy consumption on real traffic, but treats worst-case insertion cost as a secondary concern; TECache in particular reports hit rate and power savings but offers no guarantee on how long a single insertion can take before either succeeding or falling back to a full rehash. The most recent theoretical advance, bubble-up cuckoo hashing\~{}\textbackslash{}cite\{kuszmaul2025bubbling\}, achieves a strong \textbackslash{}emph\{expected\} insertion-time bound at high load, but it is an offline-insertion-time result validated against random-graph theory, not an online scheme instrumented against the specific metrics---worst-case per-insertion latency, stash utilization, and rehash frequency under a synthetic or real flow-table workload---that a line-rate SDN deployment must be evaluated on.

The gap, concretely, is this: no existing construction couples an \textbackslash{}emph\{explicitly and tightly bounded\} kick-chain length with \textbackslash{}emph\{early\} stash fallback---triggering the stash before the classical process would otherwise continue searching---while confining rehashing to the rare event of stash overflow rather than ordinary cycle detection, and then characterizes the resulting system on the metrics that matter for exact-match flow tables specifically: worst-case and average lookup latency, insertion success rate, stash utilization, and memory overhead, benchmarked against the standard baselines of chained hashing, unstashed 2-way cuckoo hashing, and linear probing on real or synthetic microflow traces.

This paper closes that gap by treating the kick-chain bound itself as the primary design variable rather than an afterthought. The proposed construction combines \$d\$-ary candidate slot selection\~{}\textbackslash{}cite\{fountoulakis2012sharp, kuszmaul2025bubbling\} with a small Kirsch-style stash\~{}\textbackslash{}cite\{kirsch2009stash\} whose sizing is informed by Noble's tightened, general-stash bounds\~{}\textbackslash{}cite\{noble2021explicit\} and corrected against the Minaud--Papamanthou erratum\~{}\textbackslash{}cite\{minaud2023generalized\}, but departs from all of these in one specific respect: the eviction procedure is capped at a deliberately short \textbackslash{}texttt\{MaxKicks\} bound, well below what a cycle-detection-only policy would tolerate, so that displaced keys are routed to the stash \textbackslash{}emph\{proactively\} rather than only after the classical process has exhausted its search. Rehashing is triggered exclusively by stash overflow and is performed incrementally---reinserting only the stash's contents---rather than as a full-table rebuild, directly targeting the catastrophic-cost scenario that has motivated every stash-based scheme since Kirsch et al.\~{}\textbackslash{}cite\{kirsch2009stash\} without adopting their assumption of an effectively unbounded search process beforehand. A load-factor monitor, sized using the exact thresholds established by Fountoulakis and Panagiotou\~{}\textbackslash{}cite\{fountoulakis2012sharp\}, resizes the table proactively rather than reactively, keeping the stash-overflow event itself rare by construction.

The resulting scheme is positioned deliberately alongside, rather than in competition with, the closest systems-level precedent: TECache\~{}\textbackslash{}cite\{xiong2025tecache\} demonstrates that a cuckoo-style cache can sustain high hit rates and energy savings for SDN TCAM flow tables, but offers no bounded worst-case insertion guarantee; this work supplies exactly that guarantee, evaluated using the same class of real and synthetic exact-match traffic that TECache and Cuckoo Cache\~{}\textbackslash{}cite\{pontarelli2016cuckoocache\} use to validate throughput, but reoriented around the central empirical question a bounded-displacement design raises: how short can the kick-chain bound be made before the resulting rise in stash utilization and rehash frequency becomes unacceptable, at a given table load factor. Answering that question, with the explicit theoretical grounding provided by\~{}\textbackslash{}cite\{noble2021explicit\} and the corrected bounds of\~{}\textbackslash{}cite\{minaud2023generalized\}, is the specific contribution this paper makes to the cuckoo hashing literature surveyed above.

\section{References}
\bibliographystyle{ieeetr}
\bibliography{references}

\end{document}
